% 原书术语索引与记号索引的完整译文；不对索引本身使用 Term 宏。
% link_terms.py 应跳过此文件，避免把索引本身列入正文术语用例。
\SGAfrontchapter{原书术语索引（中文译文）}
\SourcePage{339}
原索引按法文词头排列，以下保留其词头与排列次序，并附中文译名。右栏保留原重排本的印刷页码，括号中的链接通往译文里对应的原 PDF 页位置；两者相差 $16$ 页。词头中的“分别地”对应原文的 \textit{resp.}。
\par\medskip
\begingroup\small\setlength{\tabcolsep}{4pt}
\begin{longtable}{@{}>{\raggedright\arraybackslash}p{.71\textwidth}>{\raggedright\arraybackslash}p{.25\textwidth}@{}}
\toprule\textbf{中文／法文词头}&\textbf{原书页码／来源}\\\midrule\endfirsthead
\toprule\textbf{中文／法文词头}&\textbf{原书页码／来源}\\\midrule\endhead
\bottomrule\endfoot
\textbf{Abhyankar 引理}\newline {\itshape Abhyankar (lemme d’)} & 214（\SGAsourceRef{230}）\tabularnewline\addlinespace[.45em]
\textbf{绝对 Abhyankar 引理}\newline {\itshape Abhyankar absolu (lemme d’)} & 315（\SGAsourceRef{331}）\tabularnewline\addlinespace[.45em]
\textbf{相对 Abhyankar 引理}\newline {\itshape Abhyankar relatif (lemme d’)} & 317（\SGAsourceRef{333}）\tabularnewline\addlinespace[.45em]
\textbf{解析整}\newline {\itshape analytiquement intègre} & 28（\SGAsourceRef{44}）\tabularnewline\addlinespace[.45em]
\textbf{解析约化}\newline {\itshape analytiquement réduit} & 14（\SGAsourceRef{30}）\tabularnewline\addlinespace[.45em]
\textbf{正则增广环}\newline {\itshape anneau augmenté régulier} & 39（\SGAsourceRef{55}）\tabularnewline\addlinespace[.45em]
\textbf{Artin–Schreier 正合列}\newline {\itshape Artin-Schreier (suite exacte d’)} & 236（\SGAsourceRef{252}）\tabularnewline\addlinespace[.45em]
\textbf{自对偶性}\newline {\itshape autodualité} & 147（\SGAsourceRef{163}）\tabularnewline\addlinespace[.45em]
\textbf{双纤维范畴}\newline {\itshape bi-fibrée (catégorie)} & 143（\SGAsourceRef{159}）\tabularnewline\addlinespace[.45em]
\textbf{笛卡尔函子}\newline {\itshape cartésien (foncteur)} & 131（\SGAsourceRef{147}）\tabularnewline\addlinespace[.45em]
\textbf{笛卡尔态射}\newline {\itshape cartésien (morphisme)} & 130（\SGAsourceRef{146}）\tabularnewline\addlinespace[.45em]
\textbf{带劈分的范畴（分别地，带正规化劈分的范畴）}\newline {\itshape catégorie clivée (resp. clivée normalisée)} & 136（\SGAsourceRef{152}）\tabularnewline\addlinespace[.45em]
\textbf{纤维范畴（范畴的纤维）}\newline {\itshape catégorie-fibre} & 128（\SGAsourceRef{144}）\tabularnewline\addlinespace[.45em]
\textbf{叠}\newline {\itshape champ} & 260（\SGAsourceRef{276}）\tabularnewline\addlinespace[.45em]
\textbf{1-可构造叠（分别地，可构造叠）}\newline {\itshape champ 1-constructible (resp. constructible)} & 260（\SGAsourceRef{276}）\tabularnewline\addlinespace[.45em]
\textbf{归纳有限叠（分别地，归纳 $\ell$-有限叠）}\newline {\itshape champ ind-fini (resp. ind-$\ell$-fini)} & 260（\SGAsourceRef{276}）\tabularnewline\addlinespace[.45em]
\textbf{基变换函子}\newline {\itshape changement de base (foncteur)} & 125（\SGAsourceRef{141}）\tabularnewline\addlinespace[.45em]
\textbf{劈分}\newline {\itshape clivage} & 136（\SGAsourceRef{152}）\tabularnewline\addlinespace[.45em]
\textbf{正规化劈分}\newline {\itshape clivage normalisé} & 136（\SGAsourceRef{152}）\tabularnewline\addlinespace[.45em]
\textbf{余纤维范畴}\newline {\itshape co-fibrée (catégorie)} & 143（\SGAsourceRef{159}）\tabularnewline\addlinespace[.45em]
\textbf{Cohen–Macaulay 概形}\newline {\itshape Cohen-Macaulay (schéma de)} & 29（\SGAsourceRef{45}）\newline 43（\SGAsourceRef{59}）\tabularnewline\addlinespace[.45em]
\textbf{上同调固有}\newline {\itshape cohomologiquement propre} & 261（\SGAsourceRef{277}）\tabularnewline\addlinespace[.45em]
\textbf{完全分解的对象}\newline {\itshape complètement décomposé (objet)} & 113（\SGAsourceRef{129}）\tabularnewline\addlinespace[.45em]
\textbf{重数为 $1$ 的分支}\newline {\itshape composante de multiplicité 1} & 29（\SGAsourceRef{45}）\tabularnewline\addlinespace[.45em]
\textbf{范畴中的连通对象}\newline {\itshape connexe (objet d’une catégorie)} & 99（\SGAsourceRef{115}）\tabularnewline\addlinespace[.45em]
\textbf{余法层}\newline {\itshape conormal (faisceau)} & 33（\SGAsourceRef{49}）\tabularnewline\addlinespace[.45em]
\textbf{可构造叠}\newline {\itshape constructible (champ)} & 260（\SGAsourceRef{276}）\tabularnewline\addlinespace[.45em]
\textbf{可构造子集}\newline {\itshape constructible (partie)} & 82（\SGAsourceRef{98}）\tabularnewline\addlinespace[.45em]
\textbf{光滑性的 Jacobi 判据}\newline {\itshape critère jacobien de lissité} & 37（\SGAsourceRef{53}）\tabularnewline\addlinespace[.45em]
\textbf{正规交叉除子}\newline {\itshape croisements normaux (diviseur à)} & 277（\SGAsourceRef{293}）\tabularnewline\addlinespace[.45em]
\textbf{完全分解的对象}\newline {\itshape décomposé (objet complètement)} & 113（\SGAsourceRef{129}）\tabularnewline\addlinespace[.45em]
\textbf{分解群}\newline {\itshape décomposition (groupe de)} & 92（\SGAsourceRef{108}）\tabularnewline\addlinespace[.45em]
\textbf{导子层}\newline {\itshape dérivations (faisceau des)} & 61（\SGAsourceRef{77}）\tabularnewline\addlinespace[.45em]
\textbf{下降态射}\newline {\itshape descente (morphisme de)} & 165（\SGAsourceRef{181}）\tabularnewline\addlinespace[.45em]
\textbf{模同态的下降}\newline {\itshape descente d’homomorphismes de modules} & 154（\SGAsourceRef{170}）\tabularnewline\addlinespace[.45em]
\textbf{模的下降}\newline {\itshape descente de modules} & 154（\SGAsourceRef{170}）\tabularnewline\addlinespace[.45em]
\textbf{商模的下降}\newline {\itshape descente de modules quotients} & 156（\SGAsourceRef{172}）\tabularnewline\addlinespace[.45em]
\textbf{模的性质的下降}\newline {\itshape descente de propriétés de modules} & 157（\SGAsourceRef{173}）\tabularnewline\addlinespace[.45em]
\textbf{模截面的下降}\newline {\itshape descente de sections de modules} & 156（\SGAsourceRef{172}）\tabularnewline\addlinespace[.45em]
\textbf{开子集（分别地，闭子集）的下降}\newline {\itshape descente des parties ouvertes (resp. fermées)} & 161（\SGAsourceRef{177}）\tabularnewline\addlinespace[.45em]
\textbf{闭子预概形的下降}\newline {\itshape descente des sous-préschémas fermés} & 157（\SGAsourceRef{173}）\tabularnewline\addlinespace[.45em]
\textbf{有效下降态射}\newline {\itshape descente effective (morphisme de)} & 153（\SGAsourceRef{169}）\tabularnewline\addlinespace[.45em]
\textbf{不同理想}\newline {\itshape différente (idéal)} & 3（\SGAsourceRef{19}）\tabularnewline\addlinespace[.45em]
\textbf{微分光滑}\newline {\itshape différentiablement lisse} & 41（\SGAsourceRef{57}）\tabularnewline\addlinespace[.45em]
\textbf{相对维数}\newline {\itshape dimension relative} & 50（\SGAsourceRef{66}）\tabularnewline\addlinespace[.45em]
\textbf{正规交叉除子}\newline {\itshape diviseur à croisements normaux} & 277（\SGAsourceRef{293}）\tabularnewline\addlinespace[.45em]
\textbf{有效下降数据}\newline {\itshape donnée de descente effective} & 154（\SGAsourceRef{170}）\tabularnewline\addlinespace[.45em]
\textbf{诱导的下降数据}\newline {\itshape donnée de descente induite} & 171（\SGAsourceRef{187}）\tabularnewline\addlinespace[.45em]
\textbf{叠的温和分歧元素}\newline {\itshape élément modérément ramifié d’un champ} & 281（\SGAsourceRef{297}）\tabularnewline\addlinespace[.45em]
\textbf{有效满态射}\newline {\itshape épimorphisme effectif} & 164（\SGAsourceRef{180}）\tabularnewline\addlinespace[.45em]
\textbf{等维态射}\newline {\itshape équidimensionnel (morphisme)} & 27（\SGAsourceRef{43}）\tabularnewline\addlinespace[.45em]
\textbf{范畴等价}\newline {\itshape équivalence de catégories} & 121（\SGAsourceRef{137}）\tabularnewline\addlinespace[.45em]
\textbf{相伴解析空间}\newline {\itshape espace analytique associé} & 239（\SGAsourceRef{255}）\tabularnewline\addlinespace[.45em]
\textbf{本质平展}\newline {\itshape essentiellement étale} & 50（\SGAsourceRef{66}）\tabularnewline\addlinespace[.45em]
\textbf{本质光滑}\newline {\itshape essentiellement lisse} & 26（\SGAsourceRef{42}）\tabularnewline\addlinespace[.45em]
\textbf{本质非分歧（分别地，本质净）}\newline {\itshape essentiellement non ramifié (resp. net)} & 50（\SGAsourceRef{66}）\tabularnewline\addlinespace[.45em]
\textbf{平展态射、平展代数}\newline {\itshape étale (morphisme, algèbre)} & 4（\SGAsourceRef{20}）\newline 177（\SGAsourceRef{193}）\tabularnewline\addlinespace[.45em]
\textbf{平展覆盖}\newline {\itshape étale (revêtement)} & 5（\SGAsourceRef{21}）\tabularnewline\addlinespace[.45em]
\textbf{余法层}\newline {\itshape faisceau conormal} & 33（\SGAsourceRef{49}）\tabularnewline\addlinespace[.45em]
\textbf{导子层}\newline {\itshape faisceau des dérivations} & 61（\SGAsourceRef{77}）\tabularnewline\addlinespace[.45em]
\textbf{极大子凝集组成的层}\newline {\itshape faisceau des sous-gerbes maximales} & 260（\SGAsourceRef{276}）\tabularnewline\addlinespace[.45em]
\textbf{切层}\newline {\itshape faisceau tangent} & 61（\SGAsourceRef{77}）\tabularnewline\addlinespace[.45em]
\textbf{纤维化函子}\newline {\itshape fibrant (foncteur)} & 132（\SGAsourceRef{148}）\tabularnewline\addlinespace[.45em]
\textbf{纤维范畴}\newline {\itshape fibrée (catégorie)} & 132（\SGAsourceRef{148}）\tabularnewline\addlinespace[.45em]
\textbf{群胚纤维范畴}\newline {\itshape fibrée en groupoïdes} & 132（\SGAsourceRef{148}）\tabularnewline\addlinespace[.45em]
\end{longtable}
\SourcePage{340}
\begin{longtable}{@{}>{\raggedright\arraybackslash}p{.71\textwidth}>{\raggedright\arraybackslash}p{.25\textwidth}@{}}
\toprule\textbf{中文／法文词头}&\textbf{原书页码／来源}\\\midrule\endfirsthead
\toprule\textbf{中文／法文词头}&\textbf{原书页码／来源}\\\midrule\endhead
\bottomrule\endfoot
\textbf{忠实函子}\newline {\itshape fidèle (foncteur)} & 121（\SGAsourceRef{137}）\tabularnewline\addlinespace[.45em]
\textbf{忠实平坦模}\newline {\itshape fidèlement plat (module)} & 74（\SGAsourceRef{90}）\tabularnewline\addlinespace[.45em]
\textbf{忠实平坦态射}\newline {\itshape fidèlement plat (morphisme)} & 153（\SGAsourceRef{169}）\tabularnewline\addlinespace[.45em]
\textbf{基变换函子}\newline {\itshape foncteur changement de base} & 125（\SGAsourceRef{141}）\tabularnewline\addlinespace[.45em]
\textbf{纤维函子}\newline {\itshape foncteur fibre} & 104（\SGAsourceRef{120}）\tabularnewline\addlinespace[.45em]
\textbf{基本函子}\newline {\itshape foncteur fondamental} & 104（\SGAsourceRef{120}）\tabularnewline\addlinespace[.45em]
\textbf{形式平展}\newline {\itshape formellement étale} & 50（\SGAsourceRef{66}）\tabularnewline\addlinespace[.45em]
\textbf{形式光滑同态、形式光滑代数}\newline {\itshape formellement lisse (homomorphisme, algèbre)} & 49（\SGAsourceRef{65}）\tabularnewline\addlinespace[.45em]
\textbf{形式非分歧（分别地，形式净）}\newline {\itshape formellement non ramifié (resp. net)} & 50（\SGAsourceRef{66}）\tabularnewline\addlinespace[.45em]
\textbf{形式主齐性层}\newline {\itshape formellement principal homogène (faisceau)} & 59（\SGAsourceRef{75}）\tabularnewline\addlinespace[.45em]
\textbf{形式主齐性对象}\newline {\itshape formellement principal homogène (objet)} & 109（\SGAsourceRef{125}）\tabularnewline\addlinespace[.45em]
\textbf{形式主齐性预概形}\newline {\itshape formellement principal homogène (préschéma)} & 94（\SGAsourceRef{110}）\newline 227（\SGAsourceRef{243}）\tabularnewline\addlinespace[.45em]
\textbf{Galois 对象}\newline {\itshape galoisien (objet)} & 100（\SGAsourceRef{116}）\tabularnewline\addlinespace[.45em]
\textbf{Galois 范畴}\newline {\itshape galoisienne (catégorie)} & 104（\SGAsourceRef{120}）\tabularnewline\addlinespace[.45em]
\textbf{几何连通}\newline {\itshape géométriquement connexe} & 182（\SGAsourceRef{198}）\tabularnewline\addlinespace[.45em]
\textbf{几何平凡覆盖}\newline {\itshape géométriquement trivial (revêtement)} & 196（\SGAsourceRef{212}）\tabularnewline\addlinespace[.45em]
\textbf{几何单分支}\newline {\itshape géométriquement unibranche} & 22（\SGAsourceRef{38}）\tabularnewline\addlinespace[.45em]
\textbf{凝集}\newline {\itshape gerbe} & 260（\SGAsourceRef{276}）\tabularnewline\addlinespace[.45em]
\textbf{平凡凝集}\newline {\itshape gerbe triviale} & 260（\SGAsourceRef{276}）\tabularnewline\addlinespace[.45em]
\textbf{Grauert–Remmert 定理}\newline {\itshape Grauert-Remmert (théorème de)} & 255（\SGAsourceRef{271}）\tabularnewline\addlinespace[.45em]
\textbf{惯性群}\newline {\itshape groupe d’inertie} & 92（\SGAsourceRef{108}）\tabularnewline\addlinespace[.45em]
\textbf{分解群}\newline {\itshape groupe de décomposition} & 92（\SGAsourceRef{108}）\tabularnewline\addlinespace[.45em]
\textbf{主覆盖的 Galois 群}\newline {\itshape groupe de Galois d’un revêtement principal} & 96（\SGAsourceRef{112}）\tabularnewline\addlinespace[.45em]
\textbf{预概形在一点处的基本群}\newline {\itshape groupe fondamental d’un préschéma en un point} & 115（\SGAsourceRef{131}）\tabularnewline\addlinespace[.45em]
\textbf{Galois 范畴的基本群}\newline {\itshape groupe fondamental d’une catégorie galoisienne} & 101（\SGAsourceRef{117}）\tabularnewline\addlinespace[.45em]
\textbf{温和分歧基本群}\newline {\itshape groupe fondamental modérément ramifié} & 278（\SGAsourceRef{294}）\tabularnewline\addlinespace[.45em]
\textbf{基本群胚}\newline {\itshape groupoïde fondamental} & 107（\SGAsourceRef{123}）\tabularnewline\addlinespace[.45em]
\textbf{基本群的特殊化同态}\newline {\itshape homomorphisme de spécialisation pour le groupe fondamental} & 207（\SGAsourceRef{223}）\newline 290（\SGAsourceRef{306}）\tabularnewline\addlinespace[.45em]
\textbf{不同理想}\newline {\itshape idéal différente} & 3（\SGAsourceRef{19}）\tabularnewline\addlinespace[.45em]
\textbf{正则理想}\newline {\itshape idéal régulier} & 38（\SGAsourceRef{54}）\tabularnewline\addlinespace[.45em]
\textbf{正像}\newline {\itshape image directe} & 143（\SGAsourceRef{159}）\tabularnewline\addlinespace[.45em]
\textbf{叠的温和分歧正像}\newline {\itshape image directe modérément ramifiée d’un champ} & 281（\SGAsourceRef{297}）\tabularnewline\addlinespace[.45em]
\textbf{逆像}\newline {\itshape image inverse} & 130（\SGAsourceRef{146}）\tabularnewline\addlinespace[.45em]
\textbf{正则浸入}\newline {\itshape immersion régulière} & 38（\SGAsourceRef{54}）\tabularnewline\addlinespace[.45em]
\textbf{惯性群}\newline {\itshape inertie (groupe d’)} & 92（\SGAsourceRef{108}）\tabularnewline\addlinespace[.45em]
\textbf{基本群的双有理不变性}\newline {\itshape invariance birationnelle du groupe fondamental} & 214（\SGAsourceRef{230}）\tabularnewline\addlinespace[.45em]
\textbf{光滑性的 Jacobi 判据}\newline {\itshape jacobien (critère de lissité)} & 37（\SGAsourceRef{53}）\tabularnewline\addlinespace[.45em]
\textbf{Kummer 正合列}\newline {\itshape Kummer (suite exacte de)} & 233（\SGAsourceRef{249}）\tabularnewline\addlinespace[.45em]
\textbf{Künneth 公式}\newline {\itshape Künneth (formule de)} & 310（\SGAsourceRef{326}）\tabularnewline\addlinespace[.45em]
\textbf{Abhyankar 引理}\newline {\itshape lemme d’Abhyankar} & 214（\SGAsourceRef{230}）\newline 314（\SGAsourceRef{330}）\tabularnewline\addlinespace[.45em]
\textbf{形式光滑}\newline {\itshape lisse (formellement)} & 49（\SGAsourceRef{65}）\tabularnewline\addlinespace[.45em]
\textbf{光滑态射、光滑代数}\newline {\itshape lisse (morphisme, algèbre)} & 25（\SGAsourceRef{41}）\tabularnewline\addlinespace[.45em]
\textbf{光滑性的 Jacobi 判据}\newline {\itshape lissité (critère jacobien de)} & 37（\SGAsourceRef{53}）\tabularnewline\addlinespace[.45em]
\textbf{温和分歧层}\newline {\itshape modérément ramifié (faisceau)} & 277（\SGAsourceRef{293}）\newline 282（\SGAsourceRef{298}）\tabularnewline\addlinespace[.45em]
\textbf{温和分歧基本群}\newline {\itshape modérément ramifié (groupe fondamental)} & 278（\SGAsourceRef{294}）\tabularnewline\addlinespace[.45em]
\textbf{温和分歧万有覆盖}\newline {\itshape modérément ramifié (revêtement universel)} & 278（\SGAsourceRef{294}）\tabularnewline\addlinespace[.45em]
\textbf{温和分歧代数}\newline {\itshape modérément ramifiée (algèbre)} & 275（\SGAsourceRef{291}）\tabularnewline\addlinespace[.45em]
\textbf{叠的温和分歧正像}\newline {\itshape modérément ramifiée (image directe d’un champ)} & 281（\SGAsourceRef{297}）\tabularnewline\addlinespace[.45em]
\textbf{多重 Galois 范畴}\newline {\itshape multi-galoisienne (catégorie)} & 118（\SGAsourceRef{134}）\tabularnewline\addlinespace[.45em]
\textbf{分支的重数}\newline {\itshape multiplicité d’une composante} & 29（\SGAsourceRef{45}）\tabularnewline\addlinespace[.45em]
\textbf{净态射、净代数}\newline {\itshape net (morphisme, algèbre)} & 2（\SGAsourceRef{18}）\tabularnewline\addlinespace[.45em]
\textbf{净覆盖}\newline {\itshape net (revêtement)} & 5（\SGAsourceRef{21}）\tabularnewline\addlinespace[.45em]
\textbf{带点对象的中性连通分支}\newline {\itshape neutre (composante connexe d’un objet ponctué)} & 113（\SGAsourceRef{129}）\tabularnewline\addlinespace[.45em]
\textbf{非分歧态射、非分歧代数}\newline {\itshape non ramifié (morphisme, algèbre)} & 2（\SGAsourceRef{18}）\tabularnewline\addlinespace[.45em]
\textbf{正规化劈分}\newline {\itshape normalisé (clivage)} & 136（\SGAsourceRef{152}）\tabularnewline\addlinespace[.45em]
\textbf{将概形延拓为光滑概形的障碍}\newline {\itshape obstruction au prolongement d’un schéma en un schéma lisse} & 66（\SGAsourceRef{82}）\tabularnewline\addlinespace[.45em]
\textbf{$\leq n$ 阶微分算子}\newline {\itshape opérateur différentiel d’ordre $\leq n$} & 42（\SGAsourceRef{58}）\tabularnewline\addlinespace[.45em]
\textbf{万有 $\leq n$ 阶微分算子}\newline {\itshape opérateur différentiel d’ordre $\leq n$ universel} & 42（\SGAsourceRef{58}）\tabularnewline\addlinespace[.45em]
\textbf{$n$ 阶主部层、主部系}\newline {\itshape parties principales d’ordre $n$ (faisceau des, système des)} & 42（\SGAsourceRef{58}）\tabularnewline\addlinespace[.45em]
\textbf{忠实平坦态射}\newline {\itshape plat (morphisme fidèlement)} & 153（\SGAsourceRef{169}）\tabularnewline\addlinespace[.45em]
\textbf{平坦模、忠实平坦模}\newline {\itshape plat, fidèlement plat (module)} & 72（\SGAsourceRef{88}）\newline 74（\SGAsourceRef{90}）\tabularnewline\addlinespace[.45em]
\textbf{全忠实函子}\newline {\itshape pleinement fidèle (foncteur)} & 121（\SGAsourceRef{137}）\tabularnewline\addlinespace[.45em]
\textbf{带点对象}\newline {\itshape ponctué (objet)} & 112（\SGAsourceRef{128}）\tabularnewline\addlinespace[.45em]
\textbf{预纤维范畴}\newline {\itshape préfibrée (catégorie)} & 132（\SGAsourceRef{148}）\tabularnewline\addlinespace[.45em]
\textbf{阶数与剩余特征互素的层}\newline {\itshape premier aux caractéristiques résiduelles (faisceau)} & 281（\SGAsourceRef{297}）\tabularnewline\addlinespace[.45em]
\textbf{有限表示态射}\newline {\itshape présentation finie (morphisme de)} & 160（\SGAsourceRef{176}）\tabularnewline\addlinespace[.45em]
\textbf{主覆盖}\newline {\itshape principal (revêtement)} & 96（\SGAsourceRef{112}）\tabularnewline\addlinespace[.45em]
\textbf{主齐性层}\newline {\itshape principal homogène (faisceau)} & 59（\SGAsourceRef{75}）\tabularnewline\addlinespace[.45em]
\textbf{主齐性对象}\newline {\itshape principal homogène (objet)} & 109（\SGAsourceRef{125}）\tabularnewline\addlinespace[.45em]
\textbf{主齐性预概形、主齐性丛}\newline {\itshape principal homogène (préschéma, fibré)} & 95（\SGAsourceRef{111}）\tabularnewline\addlinespace[.45em]
\textbf{局部平凡的主齐性丛}\newline {\itshape principal homogène localement trivial (fibré)} & 230（\SGAsourceRef{246}）\tabularnewline\addlinespace[.45em]
\textbf{预概形的投基本群}\newline {\itshape pro-groupe fondamental d’un préschéma} & 116（\SGAsourceRef{132}）\tabularnewline\addlinespace[.45em]
\textbf{Galois 范畴的投基本群}\newline {\itshape pro-groupe fondamental d’une catégorie galoisienne} & 108（\SGAsourceRef{124}）\tabularnewline\addlinespace[.45em]
\textbf{基本投对象}\newline {\itshape pro-objet fondamental} & 104（\SGAsourceRef{120}）\tabularnewline\addlinespace[.45em]
\textbf{提升延拓定理}\newline {\itshape prolongement des relèvements (théorème de)} & 7（\SGAsourceRef{23}）\tabularnewline\addlinespace[.45em]
\textbf{上同调固有性}\newline {\itshape propreté cohomologique} & 261（\SGAsourceRef{277}）\tabularnewline\addlinespace[.45em]
\textbf{拟有限态射、拟有限代数}\newline {\itshape quasi-fini (morphisme, algèbre)} & 2（\SGAsourceRef{18}）\tabularnewline\addlinespace[.45em]
\textbf{商预概形}\newline {\itshape quotient (préschéma)} & 88（\SGAsourceRef{104}）\tabularnewline\addlinespace[.45em]
\end{longtable}
\SourcePage{341}
\begin{longtable}{@{}>{\raggedright\arraybackslash}p{.71\textwidth}>{\raggedright\arraybackslash}p{.25\textwidth}@{}}
\toprule\textbf{中文／法文词头}&\textbf{原书页码／来源}\\\midrule\endfirsthead
\toprule\textbf{中文／法文词头}&\textbf{原书页码／来源}\\\midrule\endhead
\bottomrule\endfoot
\textbf{根性态射}\newline {\itshape radiciel (morphisme)} & 5（\SGAsourceRef{21}）\tabularnewline\addlinespace[.45em]
\textbf{正则增广环}\newline {\itshape régulier (anneau augmenté)} & 39（\SGAsourceRef{55}）\tabularnewline\addlinespace[.45em]
\textbf{正则理想}\newline {\itshape régulier (idéal)} & 38（\SGAsourceRef{54}）\tabularnewline\addlinespace[.45em]
\textbf{正则生成元系}\newline {\itshape régulier (système de générateurs)} & 37（\SGAsourceRef{53}）\tabularnewline\addlinespace[.45em]
\textbf{正则浸入}\newline {\itshape régulière (immersion)} & 38（\SGAsourceRef{54}）\tabularnewline\addlinespace[.45em]
\textbf{正则浸入的预概形}\newline {\itshape régulièrement immergé (préschéma)} & 38（\SGAsourceRef{54}）\tabularnewline\addlinespace[.45em]
\textbf{平展覆盖（分别地，净覆盖）}\newline {\itshape revêtement étale (resp. net)} & 5（\SGAsourceRef{21}）\tabularnewline\addlinespace[.45em]
\textbf{解析空间的有限覆盖}\newline {\itshape revêtement fini d’un espace analytique} & 251（\SGAsourceRef{267}）\tabularnewline\addlinespace[.45em]
\textbf{主覆盖}\newline {\itshape revêtement principal} & 96（\SGAsourceRef{112}）\tabularnewline\addlinespace[.45em]
\textbf{预概形在一点处的万有覆盖}\newline {\itshape revêtement universel d’un préschéma en un point} & 115（\SGAsourceRef{131}）\tabularnewline\addlinespace[.45em]
\textbf{温和分歧万有覆盖}\newline {\itshape revêtement universel modérément ramifié} & 278（\SGAsourceRef{294}）\tabularnewline\addlinespace[.45em]
\textbf{Riemann 存在定理}\newline {\itshape Riemann (théorème d’existence de)} & 251（\SGAsourceRef{267}）\tabularnewline\addlinespace[.45em]
\textbf{分裂}\newline {\itshape scindage} & 142（\SGAsourceRef{158}）\tabularnewline\addlinespace[.45em]
\textbf{分裂范畴}\newline {\itshape scindée (catégorie)} & 142（\SGAsourceRef{158}）\tabularnewline\addlinespace[.45em]
\textbf{可分多项式}\newline {\itshape séparable (polynôme)} & 9（\SGAsourceRef{25}）\tabularnewline\addlinespace[.45em]
\textbf{Serre 判据}\newline {\itshape Serre (critère de)} & 15（\SGAsourceRef{31}）\tabularnewline\addlinespace[.45em]
\textbf{简单：“光滑”的旧称}\newline {\itshape simple : synonyme désuet de lisse} & 25（\SGAsourceRef{41}）\tabularnewline\addlinespace[.45em]
\textbf{纤维子范畴（分别地，预纤维子范畴）}\newline {\itshape sous-catégorie fibrée (resp. préfibrée)} & 132（\SGAsourceRef{148}）\tabularnewline\addlinespace[.45em]
\textbf{基本群的特殊化同态}\newline {\itshape spécialisation pour le groupe fondamental (homomorphisme de)} & 207（\SGAsourceRef{223}）\newline 290（\SGAsourceRef{306}）\tabularnewline\addlinespace[.45em]
\textbf{严格正规交叉除子}\newline {\itshape strictement à croisements normaux (diviseur)} & 277（\SGAsourceRef{293}）\tabularnewline\addlinespace[.45em]
\textbf{淹没态射}\newline {\itshape submersif (morphisme)} & 179（\SGAsourceRef{195}）\tabularnewline\addlinespace[.45em]
\textbf{截面的 $n$ 阶主部系}\newline {\itshape système des parties principales d’ordre $n$ d’une section} & 42（\SGAsourceRef{58}）\tabularnewline\addlinespace[.45em]
\textbf{基本群的局部系统}\newline {\itshape système local des groupes fondamentaux} & 116（\SGAsourceRef{132}）\tabularnewline\addlinespace[.45em]
\textbf{正则生成元系}\newline {\itshape système régulier de générateurs} & 37（\SGAsourceRef{53}）\tabularnewline\addlinespace[.45em]
\textbf{切层}\newline {\itshape tangent (faisceau)} & 61（\SGAsourceRef{77}）\tabularnewline\addlinespace[.45em]
\textbf{Riemann 存在定理}\newline {\itshape théorème d’existence de Riemann} & 251（\SGAsourceRef{267}）\tabularnewline\addlinespace[.45em]
\textbf{基本群的双有理不变性定理}\newline {\itshape théorème d’invariance birationnelle du groupe fondamental} & 214（\SGAsourceRef{230}）\tabularnewline\addlinespace[.45em]
\textbf{Grauert–Remmert 定理}\newline {\itshape théorème de Grauert-Remmert} & 255（\SGAsourceRef{271}）\tabularnewline\addlinespace[.45em]
\textbf{提升延拓定理}\newline {\itshape théorème de prolongement des relèvements} & 7（\SGAsourceRef{23}）\tabularnewline\addlinespace[.45em]
\textbf{纯性定理}\newline {\itshape théorème de pureté} & 212（\SGAsourceRef{228}）\tabularnewline\addlinespace[.45em]
\textbf{搬运态射}\newline {\itshape transport (morphisme de)} & 136（\SGAsourceRef{152}）\tabularnewline\addlinespace[.45em]
\textbf{几何单分支}\newline {\itshape unibranche (géométriquement)} & 22（\SGAsourceRef{38}）\tabularnewline\addlinespace[.45em]
\textbf{单有理簇}\newline {\itshape unirationnelle (variété)} & 219（\SGAsourceRef{235}）\tabularnewline\addlinespace[.45em]
\textbf{普遍连通：“几何连通”的同义词}\newline {\itshape universellement connexe : synonyme de géométriquement connexe} & 182（\SGAsourceRef{198}）\tabularnewline\addlinespace[.45em]
\textbf{普遍单射的模同态}\newline {\itshape universellement injectif (morphisme de modules)} & 31（\SGAsourceRef{47}）\tabularnewline\addlinespace[.45em]
\textbf{普遍开（分别地，普遍闭、普遍双连续等）}\newline {\itshape universellement ouvert (resp. universellement fermé, resp. universellement bicontinu, etc.)} & 162（\SGAsourceRef{178}）\tabularnewline\addlinespace[.45em]
\textbf{普遍淹没态射}\newline {\itshape universellement submersif (morphisme)} & 179（\SGAsourceRef{195}）\tabularnewline\addlinespace[.45em]
\textbf{$X/Y$ 的无穷小邻域}\newline {\itshape voisinage infinitésimal de $X/Y$} & 1（\SGAsourceRef{17}）\tabularnewline\addlinespace[.45em]
\end{longtable}\endgroup
\SourcePage{342}
% 原 PDF 第342页为空白。
\SGAfrontchapter{原书记号索引（中文译文）}
\SourcePage{343}
以下保留原索引的记号及排列次序。粗体与普通体、书写体与正体、上下标以及横线各按原页区分；“或”“与”分别译自原文的 \textit{ou}、\textit{et}。右栏页码与来源链接的含义同前表。
\par\medskip
\begingroup\small\setlength{\tabcolsep}{4pt}
\begin{longtable}{@{}>{\raggedright\arraybackslash}p{.71\textwidth}>{\raggedright\arraybackslash}p{.25\textwidth}@{}}
\toprule\textbf{记号}&\textbf{原书页码／来源}\\\midrule\endfirsthead
\toprule\textbf{记号}&\textbf{原书页码／来源}\\\midrule\endhead
\bottomrule\endfoot
$\Delta_{X/Y}$，或简记为 $\Delta$ & 1（\SGAsourceRef{17}）\tabularnewline\addlinespace[.45em]
$\Omega^1_{X/Y}$ & 1（\SGAsourceRef{17}）\tabularnewline\addlinespace[.45em]
$\mathscr P^n_{X/Y}$ & 1（\SGAsourceRef{17}）\newline 42（\SGAsourceRef{58}）\tabularnewline\addlinespace[.45em]
$\Delta^n_{X/Y}$ & 1（\SGAsourceRef{17}）\tabularnewline\addlinespace[.45em]
$\mathfrak d_{X/Y}$ & 3（\SGAsourceRef{19}）\tabularnewline\addlinespace[.45em]
$d^n_{X/S}$ & 42（\SGAsourceRef{58}）\tabularnewline\addlinespace[.45em]
$\mathfrak g_{X/S}$ & 61（\SGAsourceRef{77}）\tabularnewline\addlinespace[.45em]
$\mathscr C(\pi)$ & 98（\SGAsourceRef{114}）\tabularnewline\addlinespace[.45em]
$\pi$ & 101（\SGAsourceRef{117}）\tabularnewline\addlinespace[.45em]
$\operatorname{Pro}\text{-}\mathscr C(\pi)$ & 105（\SGAsourceRef{121}）\tabularnewline\addlinespace[.45em]
$\Gamma$ & 107（\SGAsourceRef{123}）\tabularnewline\addlinespace[.45em]
$\pi_1(S,a)$ & 115（\SGAsourceRef{131}）\tabularnewline\addlinespace[.45em]
$\pi_1(S;a,a')$ & 115（\SGAsourceRef{131}）\tabularnewline\addlinespace[.45em]
$\mathscr C(S)$ & 116（\SGAsourceRef{132}）\tabularnewline\addlinespace[.45em]
$\pi_1(f;a')$ & 116（\SGAsourceRef{132}）\tabularnewline\addlinespace[.45em]
$\mathbf{Ens}$ & 120（\SGAsourceRef{136}）\tabularnewline\addlinespace[.45em]
$\mathbf{Cat}$ & 120（\SGAsourceRef{136}）\tabularnewline\addlinespace[.45em]
$\operatorname{Ob}(\mathscr C)$ & 120（\SGAsourceRef{136}）\tabularnewline\addlinespace[.45em]
$\operatorname{Fl}(\mathscr C)$ & 120（\SGAsourceRef{136}）\tabularnewline\addlinespace[.45em]
$\mathbf{Hom}(\mathscr C,\mathscr C')$ & 120（\SGAsourceRef{136}）\tabularnewline\addlinespace[.45em]
$\mathscr C^\circ$ & 120（\SGAsourceRef{136}）\tabularnewline\addlinespace[.45em]
$\mathbf{Cat}_{/\mathscr E}$ & 121（\SGAsourceRef{137}）\tabularnewline\addlinespace[.45em]
$\mathbf{Hom}_{\mathscr E/-}(\mathscr F,\mathscr G)$ & 122（\SGAsourceRef{138}）\tabularnewline\addlinespace[.45em]
$v*u$ & 123（\SGAsourceRef{139}）\tabularnewline\addlinespace[.45em]
$\mathscr F\times_{\mathscr E}\mathscr G$ & 124（\SGAsourceRef{140}）\tabularnewline\addlinespace[.45em]
$\lambda^*:\mathbf{Cat}_{/\mathscr E}\to\mathbf{Cat}_{/\mathscr E'}$ & 125（\SGAsourceRef{141}）\tabularnewline\addlinespace[.45em]
$\boldsymbol\Gamma(\mathscr G/\mathscr E)$ 与 $\Gamma(\mathscr G/\mathscr E)$ & 127（\SGAsourceRef{143}）\tabularnewline\addlinespace[.45em]
$\mathscr F_S$ & 128（\SGAsourceRef{144}）\tabularnewline\addlinespace[.45em]
$f^*_{\mathscr F}(\xi)$ 或 $f^*(\xi)$ & 130（\SGAsourceRef{146}）\tabularnewline\addlinespace[.45em]
$\alpha_f(\xi)$ & 130（\SGAsourceRef{146}）\tabularnewline\addlinespace[.45em]
$\mathbf{Hom}_{\mathrm{cart}}(\mathscr F,\mathscr G)$ & 131（\SGAsourceRef{147}）\tabularnewline\addlinespace[.45em]
$\mathbf{Cat}^{\mathrm{cart}}_{/\mathscr E}$ & 131（\SGAsourceRef{147}）\tabularnewline\addlinespace[.45em]
$\mathop{\operatorname{Lim}}\limits_{\longleftarrow}\mathscr F/\mathscr E$ & 131（\SGAsourceRef{147}）\tabularnewline\addlinespace[.45em]
$f_*^{\mathscr F}$ 或 $f_*$ & 143（\SGAsourceRef{159}）\tabularnewline\addlinespace[.45em]
$\mathbf{Sch}$ & 153（\SGAsourceRef{169}）\tabularnewline\addlinespace[.45em]
$\boldsymbol\mu_{n\,S}$ & 233（\SGAsourceRef{249}）\tabularnewline\addlinespace[.45em]
$X^{\mathrm{an}}$ & 239（\SGAsourceRef{255}）\tabularnewline\addlinespace[.45em]
$f^{\mathrm{an}}$ & 239（\SGAsourceRef{255}）\tabularnewline\addlinespace[.45em]
$SF$ 或 $S(F)$ & 260（\SGAsourceRef{276}）\tabularnewline\addlinespace[.45em]
$\mathrm H^1_t(U,F)$ & 277（\SGAsourceRef{293}）\tabularnewline\addlinespace[.45em]
$\mathrm R^1_tg_*F$ & 278（\SGAsourceRef{294}）\tabularnewline\addlinespace[.45em]
$C_t((U,X)/S)$ 或 $C_t$ & 278（\SGAsourceRef{294}）\tabularnewline\addlinespace[.45em]
$\pi_1^t((U,X)/S,a)$ 或 $\pi_1^t(U,a)$ 或 $\pi_1^t(U)$ & 278（\SGAsourceRef{294}）\tabularnewline\addlinespace[.45em]
$(g_*^t\Phi)_{T'}$ & 281（\SGAsourceRef{297}）\tabularnewline\addlinespace[.45em]
$\mathrm H^0(V,C_V)^\Pi$ & 283（\SGAsourceRef{299}）\tabularnewline\addlinespace[.45em]
$\pi_1^{\mathbb L}(U,a)$ & 289（\SGAsourceRef{305}）\tabularnewline\addlinespace[.45em]
$\pi_1'(X,a)$ & 306（\SGAsourceRef{322}）\tabularnewline\addlinespace[.45em]
$\pi_1^{\mathbb L}(X/S,g,\bar s)$ 或 $\pi_1^{\mathbb L}(X/S,g)$ & 309（\SGAsourceRef{325}）\tabularnewline\addlinespace[.45em]
$\pi_1^{\mathbb L}(X_{\bar s},a)_K$ & 312（\SGAsourceRef{328}）\tabularnewline\addlinespace[.45em]
$\mathbb Z_\ell[1]$ & 316（\SGAsourceRef{332}）\tabularnewline\addlinespace[.45em]
\end{longtable}\endgroup
